\documentclass[10pt]{article}
\usepackage[margin=1in]{geometry}
\usepackage[T1]{fontenc}
\usepackage{lmodern,microtype,amsmath,amssymb,amsthm}
\theoremstyle{definition}
\newtheorem{assumption}{Assumption}
\begin{document}
\section*{High-Dimensional Constrained Quasar Minimization}
We minimize a possibly nonconvex function $f:\mathbb{R}^d\to\mathbb{R}$ over a compact convex set $X$ of diameter at most $D>0$. On $X$, its gradient is $L$-Lipschitz, with $L>0$, and $f$ is $\gamma$-quasar-convex relative to a minimizer $x^\star\in X$, with $0<\gamma\leq1$. Two oracles return unbiased stochastic gradients with variance at most $\sigma^2$, where $\sigma\geq0$, using independent binary samples, and exact projections onto $X$. For $\varepsilon>0$ and $0<\delta\leq1/4$, the goal is a measurable output $\bar{x}$ satisfying $\bar{x}\in X$ and $f(\bar{x})-f(x^\star)\leq\varepsilon$ with probability at least $1-\delta$.

The optimization problem is
\[
\min_{x\in X} f(x).
\]
We use the Euclidean inner product and norm on $\mathbb{R}^d$. A starting point $x^0\in X$ is given. The algorithm receives $d,x^0,L,D,\gamma,\sigma,\varepsilon,\delta$ and accesses the objective and feasible set through the two oracles.

The parameters satisfy
\[
\begin{aligned}
&d\in\mathbb{N},\qquad L,D,\varepsilon>0,\qquad \sigma\geq0,\\
&0<\gamma\leq1,\qquad 0<\delta\leq\tfrac14,\\
&10^8\varepsilon\leq LD^2,\\
&\sigma=0\quad\text{or}\quad64\gamma\varepsilon\leq\sigma D.
\end{aligned}
\]
To state the dimension requirement, define
\[
m=\left\lfloor
\frac{\sqrt{LD^2/\varepsilon}}{10000\gamma}
+\frac{\sigma^2D^2\log(1/\delta)}{8192\gamma^2\varepsilon^2}
\right\rfloor.
\]
The required regime is
\[
2^{34}(m+1)^5\leq d.
\]

Our assumptions:
\begin{assumption}[Feasible set and minimizer]
The set $X\subseteq\mathbb{R}^d$ is compact and convex, contains $x^0$ and a fixed point $x^\star$, and satisfies
\[
\begin{aligned}
\|x-y\|&\leq D &&\text{for all }x,y\in X,\\
f(x^\star)&\leq f(x) &&\text{for all }x\in X.
\end{aligned}
\]
\end{assumption}

\begin{assumption}[Smoothness and quasar convexity]
The function $f$ is Fr\'echet differentiable at every point of $X$. For all $x,y\in X$,
\[
\|\nabla f(x)-\nabla f(y)\|\leq L\|x-y\|.
\]
Its quasar-convexity condition is
\[
\gamma\bigl(f(x)-f(x^\star)\bigr)
\leq\langle\nabla f(x),x-x^\star\rangle,
\qquad x\in X.
\]
\end{assumption}

\begin{assumption}[Stochastic gradient oracle]
There is a probability distribution $\nu$ on $\{0,1\}$ and a jointly measurable map
$g:\mathbb{R}^d\times\{0,1\}\to\mathbb{R}^d$.
For each $x\in X$, the vector $g(x,\xi)$ is Bochner integrable and its squared error is integrable under $\xi\sim\nu$, with
\[
\begin{aligned}
\mathbb{E}_{\xi\sim\nu}[g(x,\xi)]&=\nabla f(x),\\
\mathbb{E}_{\xi\sim\nu}\bigl[\|g(x,\xi)-\nabla f(x)\|^2\bigr]
&\leq\sigma^2.
\end{aligned}
\]
The samples $\xi_0,\xi_1,\ldots$ are measurable, mutually independent, have common law $\nu$, and are independent of the algorithm's randomness. At the $k$th oracle call, a stochastic gradient query at $x$ returns $g(x,\xi_k)$.
\end{assumption}

\begin{assumption}[Projection oracle]
The projection oracle returns $\Pi_X(z)$ at $z\in\mathbb{R}^d$, where
\[
\begin{aligned}
\Pi_X(z)&\in X,\\
\langle z-\Pi_X(z),w-\Pi_X(z)\rangle&\leq0,
\qquad w\in X.
\end{aligned}
\]
\end{assumption}

Both oracles accept queries anywhere in $\mathbb{R}^d$; the gradient and moment conditions above apply on $X$. Stochastic gradient calls and projection calls are counted separately. After the allotted calls, the output $\bar{x}\in\mathbb{R}^d$ must be measurable and satisfy
\[
\mathbb{P}_{P}\!\left(
\bar{x}\in X\ \text{and}\ f(\bar{x})-f(x^\star)\leq\varepsilon
\right)\geq1-\delta,
\]
where $P$ is the joint law of the algorithm's randomness and the independent oracle samples.
\end{document}
